\documentclass[11pt]{article}
\usepackage{amsmath, amssymb}
\usepackage{geometry}
\geometry{a4paper, margin=1in}
\usepackage{hyperref}

\title{Theory of Rigid Body Motion\\and Linearized Error Dynamics}
\author{}
\date{}

\begin{document}
\maketitle

\section{Introduction}
This document outlines the theoretical foundations for the MATLAB script \path{example_torque_free_motion.m}. The script simulates the rotational dynamics of a rigid body (a CubeSat) subject to zero external torques. It propagates a nominal trajectory and compares it against slightly perturbed trajectories using both the full nonlinear equations of motion and a linearized error model. The same framework also extends naturally to the case where an external torque depends on attitude angle.

\section{Nonlinear Equations of Motion}
The rotational state of the rigid body is defined by its attitude quaternion $q$ (representing the rotation from the inertial frame to the body frame) and its angular velocity $\omega$ expressed in the body frame.

\subsection{Attitude Kinematics}
We use the scalar-first quaternion convention $q = [q_0; \vec{q}_v]^T$. The kinematic differential equation governing the evolution of the attitude quaternion is:
\begin{equation}
    \dot{q} = \frac{1}{2} q \otimes \begin{bmatrix} 0 \\ \omega \end{bmatrix}
\end{equation}
where $\otimes$ denotes quaternion multiplication.

\subsection{Rigid Body Dynamics (Euler's Equations)}
The rotational dynamics of a rigid body are governed by Euler's equation. In the most general case considered here,
\begin{equation}
    J \dot{\omega} + \omega \times (J\omega) = \tau_{ext}(q, \omega, t)
\end{equation}
where $J$ is the principal moment of inertia matrix of the body and $\tau_{ext}$ is the applied external torque, expressed in the body frame.

For the torque-free script, $\tau_{ext} = 0$, so the equation reduces to:
\begin{equation}
    J \dot{\omega} + \omega \times (J\omega) = 0 \quad \implies \quad \dot{\omega} = -J^{-1} (\omega \times (J\omega))
\end{equation}

If a torque depends on attitude angle, then the nonlinear dynamics simply keep the same structure with a nonzero right-hand side. Typical examples are gravity-gradient torques, magnetic torques, spring-like restoring torques, or control torques generated from an attitude error law.

\section{Attitude Error Representation}
To study perturbations, we define a nominal trajectory $(q_{nom}(t), \omega_{nom}(t))$ and a perturbed trajectory $(q_{pert}(t), \omega_{pert}(t))$. 

The attitude error between the nominal body frame and the perturbed body frame is represented by the error quaternion $q_{err}$, defined via the multiplicative attitude error:
\begin{equation}
    q_{err} = q_{nom}^{-1} \otimes q_{pert}
\end{equation}
For small angular deviations, the error quaternion can be approximated as:
\begin{equation}
    q_{err} \approx \begin{bmatrix} 1 \\ \frac{1}{2} \delta\theta \end{bmatrix}
\end{equation}
where $\delta\theta \in \mathbb{R}^3$ is the small-angle error vector (Euler angles) relating the nominal and perturbed body frames. Therefore, the attitude error vector can be extracted as $\delta\theta \approx 2 \vec{q}_{err, v}$.

The angular rate error is simply defined as the difference between the perturbed and nominal angular rates:
\begin{equation}
    \delta\omega = \omega_{pert} - \omega_{nom}
\end{equation}

\section{Linearized Error Dynamics}
To analyze stability and design linear controllers, we linearize the nonlinear kinematics and dynamics about the time-varying nominal trajectory $(\omega_{nom}(t))$. 

\subsection{Linearized Kinematics}
Taking the derivative of the attitude error vector $\delta\theta$ expressed in the rotating nominal body frame yields an extra Coriolis-like term:
\begin{equation}
    \delta\dot{\theta} = \delta\omega - \omega_{nom} \times \delta\theta
\end{equation}
This term accounts for the relative rotation of the reference frame in which the attitude error is measured.

\subsection{Linearized Dynamics}
For the torque-free case, we substitute $\omega = \omega_{nom} + \delta\omega$ into Euler's equation:
\begin{equation}
    J (\dot{\omega}_{nom} + \delta\dot{\omega}) = -(\omega_{nom} + \delta\omega) \times J(\omega_{nom} + \delta\omega)
\end{equation}
Expanding the cross products and noting that $J\dot{\omega}_{nom} = -\omega_{nom} \times J\omega_{nom}$:
\begin{equation}
    J\delta\dot{\omega} \approx - \delta\omega \times (J\omega_{nom}) - \omega_{nom} \times (J\delta\omega)
\end{equation}
where we have dropped the second-order term $\delta\omega \times (J\delta\omega)$.
Using the anti-commutative property of the cross product ($a \times b = -b \times a$), we can rewrite $-\delta\omega \times (J\omega_{nom})$ as $(J\omega_{nom}) \times \delta\omega$. This yields the linearized dynamics implemented in the script:
\begin{equation}
    \delta\dot{\omega} = J^{-1} \Big[ (J\omega_{nom}) \times \delta\omega - \omega_{nom} \times (J\delta\omega) \Big]
\end{equation}
Using the skew-symmetric matrix operator $[\cdot]^\times$, this can be written as:
\begin{equation}
    \delta\dot{\omega} = J^{-1} \Big[ [J\omega_{nom}]^\times \delta\omega - [\omega_{nom}]^\times J \delta\omega \Big]
\end{equation}

\subsection{If the Torque Depends on Angle}
Now suppose the external torque depends on attitude, or more generally on attitude, rate, and time:
\begin{equation}
    J \dot{\omega} + \omega \times (J\omega) = \tau_{ext}(q, \omega, t)
\end{equation}
Let the nominal trajectory satisfy
\begin{equation}
    J \dot{\omega}_{nom} + \omega_{nom} \times (J\omega_{nom}) = \tau_{nom}
\end{equation}
with $\tau_{nom} = \tau_{ext}(q_{nom}, \omega_{nom}, t)$. Linearizing the torque term gives
\begin{equation}
    \delta\tau \approx K_{\theta}(t) \, \delta\theta + K_{\omega}(t) \, \delta\omega
\end{equation}
where the Jacobian matrices are local sensitivity matrices of torque with respect to attitude error and rate error along the nominal trajectory.

The linearized rotational dynamics become
\begin{equation}
    J\delta\dot{\omega} = [J\omega_{nom}]^\times \delta\omega - [\omega_{nom}]^\times J \delta\omega + \delta\tau
\end{equation}
or, after substituting the local torque model,
\begin{equation}
    \delta\dot{\omega} = J^{-1} \Big( [J\omega_{nom}]^\times \delta\omega - [\omega_{nom}]^\times J \delta\omega + K_{\theta} \, \delta\theta + K_{\omega} \, \delta\omega \Big)
\end{equation}
while the kinematics remain
\begin{equation}
    \delta\dot{\theta} = \delta\omega - \omega_{nom} \times \delta\theta
\end{equation}

Therefore, an angle-dependent torque appears in the linearized model as an additional stiffness-like term $K_{\theta} \, \delta\theta$.

\subsection{Important Special Case: Restoring Torque}
If the nominal motion is an equilibrium with $\omega_{nom} = 0$ and the torque is a restoring torque near that equilibrium, then a common local model is
\begin{equation}
    \tau_{ext}(\theta) \approx -K \, \delta\theta
\end{equation}
where $K$ is an attitude stiffness matrix. The linearized equations reduce to
\begin{align}
    \delta\dot{\theta} &= \delta\omega \\
    J\delta\dot{\omega} &= -K \, \delta\theta
\end{align}
Combining the two gives the familiar second-order small-angle system
\begin{equation}
    J \delta\ddot{\theta} + K \, \delta\theta = 0
\end{equation}
which is the rotational analog of a mass-spring oscillator.

If damping or active rate feedback is present, for example $\tau_{ext} \approx -K \, \delta\theta - C \, \delta\omega$, then
\begin{equation}
    J \delta\ddot{\theta} + C \, \delta\dot{\theta} + K \, \delta\theta = 0
\end{equation}
which is the rotational analog of a damped mass-spring system.

\subsection{Nonlinear Angle-Dependent Torques}
If the exact torque is nonlinear in angle, for example
\begin{equation}
    \tau_{ext}(\theta) = -k \sin \theta
\end{equation}
then the full nonlinear propagation must retain that exact torque law. The small-angle linearization replaces $\sin \theta$ with $\theta$, giving
\begin{equation}
    \tau_{ext}(\theta) \approx -k\theta
\end{equation}
which is accurate only while the attitude error remains small.

\section{Conservation Laws}
In the absence of external torques, two key quantities remain strictly conserved throughout the motion:
\begin{enumerate}
    \item \textbf{Angular Momentum:} The angular momentum vector $\vec{H} = J\omega$ is constant in the inertial frame. Consequently, its magnitude in the body frame, $\|\vec{H}\| = \|J\omega\|$, remains constant.
    \item \textbf{Rotational Kinetic Energy:} The rotational kinetic energy $T_{rot} = \frac{1}{2} \omega^T J \omega$ is constant.
\end{enumerate}
The script verifies the numerical accuracy of the propagation by checking the drift of these two conserved quantities over time.

With a nonzero external torque, these conservation laws change:
\begin{itemize}
    \item Angular momentum is generally no longer constant, because the external torque changes it according to $\dot{\vec{H}} = \tau_{ext}$ in the inertial frame.
    \item Rotational kinetic energy is generally no longer constant either, because the torque can do work on the body.
    \item If the torque is conservative and derives from a potential $U(\theta)$, then total mechanical energy,
    \begin{equation}
        E = \frac{1}{2} \omega^T J \omega + U(\theta)
    \end{equation}
    is conserved in the absence of damping and explicit time dependence.
\end{itemize}

So, for the current script, $\tau_{ext} = 0$ and the motion is torque-free. If you add a torque that depends on angle, the nonlinear dynamics gain a forcing term and the linearized error model gains the additional stiffness-like contribution described above.

\end{document}
